\section{Auslegung der Bremse}
Die Auslegung der Bremse erfolgt gemäß VDI 2241-1 \cite{VDI2241-1}. \\
Zentral für die Auslegung der Bremse ist das Brems-/Kupplungsmoment $M_\mathrm{K}$. Zur Berechnung dieser Größe sind zunächst die Massenträgheiten zu bestimmen.  

\subsection{Trägheiten}
Grundlegend wird zwischen dem Trägheitsmoment des Antriebs $J_\mathrm{A}$ und dem des Abtriebs/Last $J_\mathrm{L}$ unterschieden. Als Last wird dabei alles nach der Bremse gerechnet; als Antrieb alles davor. Besonders beachtet werden muss hierbei die "Reduktion von Massen und Trägheiten (vgl. \cite[S. 268f.]{Sauer}). Da kleinere Trägheiten (z.B. von Sicherungselementen) vernachlässigt werden können, werden in den folgenden Rechnungen nur Wellen und Zahnräder berücksichtigt. 
Die Trägheiten der Wellen und Zahnräder sind, durch die Auslegung in KISSsoft und die Modellierung in Solid Edge, bekannt als: 

\begin{table}[H]
    \centering
    \caption{Gegebene Trägheiten}
    \begin{tabular}{cccc}
        \toprule
         Bauteil&  Formelzeichen& Größe [$\mathrm{kg\cdot mm^2}$]& Bemerkung\\
         \midrule
         Motor& $J_\mathrm{M}$& $316$& \\
         Antriebswelle& $J_\mathrm{AW}$& $52,693$ & \\
         Zwischenwelle 1& $J_\mathrm{ZW1}$& $150,739$& \\
         Zwischenwelle 2& $J_\mathrm{ZW2}$& $515,298$& \\
         Riemenwelle& $J_\mathrm{RW}$& $5268.556$& \\
         Sonnenwelle& $J_\mathrm{SW}$& $841,285$& \\
         Planetenwelle& $J_\mathrm{PW}$& $6\times10^{3}$& \\
         Steg& $J_\mathrm{ST}$& $3,74\times10^{3}$& inkl. Anschlussplatte\\
         \midrule
         Zahnrad 1& $J_\mathrm{Z1}$& $16,81$& \\
         Zahnrad 2& $J_\mathrm{Z2}$& $3874$& \\
         Zahnrad 3& $J_\mathrm{Z3}$& $102,7$& \\
         Zahnrad 4& $J_\mathrm{Z4}$& $23080$& \\
         Zahnrad 5& $J_\mathrm{Z5}$& $1600$& \\
         Zahnrad 6& $J_\mathrm{Z6}$& $24020$& \\
         Planetenrad& $J_\mathrm{ZP}$& $2,4\times 10^{4}$& \\
         Riemenrad& $J_\mathrm{RM}$& $3,19\times 10^{5}$& \\
         \bottomrule
    \end{tabular}
    \label{Arm Gegeben}
\end{table}\\

\subsubsection{Reduktion der Massen und Trägheiten}
Wie bereits erwähnt, müssen Trägheiten und Massen auf die betrachtete Welle (jene der Bremse) reduziert werden. Dabei werden Trägheiten zusammengefasst und mit ihren Übersetzungen verrechnet; rotierende Massen werden in Trägheiten umgewandelt. Für die Reduktion einer abtriebsseitigen (nachgeschalteten) Welle gilt \cite[S. 269, Gl. 4.23]{Sauer}: 
\begin{equation}
    J_\mathrm{ab,red} = \frac{1}{i^2} J_\mathrm{ab}
\end{equation}
Und für antriebsseitige Wellen \cite[S. 270, Gl. 4.25]{Sauer}: 
\begin{equation}
    J_\mathrm{an,red} = i^2 J_\mathrm{an}
\end{equation}
Rotierende abtriebsseitige Massen werden folgendermaßen in Trägheiten umgewandelt \cite[S. 270, Gl. 4.28]{Sauer}: 
\begin{equation}
    J_\mathrm{ab,red} = m_\mathrm{ab} \left( \frac{v_\mathrm{ab}}{\omega_\mathrm{an}} \right)^2
\end{equation}\\
Die Welle der Bremse, also die Welle auf die alle anderen reduziert werden, wird selbst nicht reduziert (vgl. \cite[s. 273, Gl. 4.41]{Sauer}), kann bei An- oder Abtrieb verrechnet werden. 
Der Übersichtlichkeit halber werden zunächst die einzelnen Getriebestufen sinnvolle Systeme zusammengefasst:
\begin{equation}
    J_\mathrm{An} = J_\mathrm{M} + J_\mathrm{AW} + J_\mathrm{Z1}
\end{equation}
\begin{equation}
    J_1 = J_\mathrm{ZW1} + J_\mathrm{Z2} + J_\mathrm{Z3}
\end{equation}
\begin{equation}
    J_2 = J_\mathrm{ZW2} + J_\mathrm{Z4} + J_\mathrm{Z5}
\end{equation}
\begin{equation}
    J_R = J_\mathrm{RW} + J_\mathrm{Z6} + J_\mathrm{RM}
\end{equation}
\begin{equation}
    J_\mathrm{S} = J_\mathrm{SW} + J_\mathrm{RM}
\end{equation}
\begin{equation}
    J_\mathrm{Ab} = J_\mathrm{ST} + 3J_\mathrm{PW} + 3J_\mathrm{ZP}
\end{equation}
\noindent
Für die reduzierte Trägheit der Antriebsseite $J_\mathrm{an,red}$ folgt: 
\begin{equation}
    J_\mathrm{an,red} = J_\mathrm{An}\cdot i_\mathrm{12}^2 = 11082,191\,\mathrm{kg\cdot mm^2}
\end{equation}
Für die Trägheit der Lastseite $J_\mathrm{ab,red}$:
\begin{equation}
    J_\mathrm{ab,red} = J_\mathrm{1} + J_\mathrm{2}\frac{1}{i_\mathrm{34}^2} + (J_\mathrm{R}+J_\mathrm{S})\frac{1}{i_\mathrm{34}^2 \cdot i_\mathrm{56}^2} + J_\mathrm{Ab}\frac{1}{i_\mathrm{34}^2 \cdot i_\mathrm{56}^2 \cdot i_\mathrm{PL}^2} + m_\mathrm{ab} \left( \frac{v_\mathrm{ab}}{\omega_\mathrm{an}} \right)^2
\end{equation}
\begin{equation*}
    J_\mathrm{ab,red} = 411926,663\,\mathrm{kg\cdot mm^2}
\end{equation*}
Die gesamte reduzierte Trägheit $J_\mathrm{ges,red}$ ist also:
\begin{equation} \label{Jges}
    J_\mathrm{ges,red} = J_\mathrm{an,red} + J_\mathrm{ab,red} = 0,423\,\mathrm{kg\cdot m^2}
\end{equation}

\subsection{Kennmoment $M_\mathrm{K}$}
Das Kennmoment $M_\mathrm{K}$ muss aufgebracht werden um den Kupplungs- bzw. Bremsvorgang durchführen zu können. Da davon auszugehen ist, dass der Motor im Falle einer Notbremsung sofort aufhört zu arbeiten, muss kein Antriebsmoment $M_\mathrm{A}$ berücksichtigt werden. Als Lastmoment $M_\mathrm{L}$ wirkt der Arm am oberen Getriebeausgang. Es ist außerdem keine Anstiegszeit $t_\mathrm{12}$ zu berücksichtigen. Das Kennmoment kann demnach entsprechend \cite[S. 13, Gl. 1]{VDI2241-1} berechnet werden:
\begin{equation} \label{MK}
    M_\mathrm{K} = M_\mathrm{L} + \frac{J_\mathrm{ges,red}\cdot |\Delta \omega|}{t_3}
\end{equation}
Das Lastmoment $M_\mathrm{L}$ entsteht durch das "Fallen" des Arms in einer Notstopsituation. In diesem Fall wirkt nur die Erdbeschleunigung auf die Masse $m_\mathrm{Last}$ des Arms ein und das Lastmoment ergibt sich zu:
\begin{equation} \label{ML}
    M_\mathrm{L} = m_\mathrm{Last} \cdot g \cdot l_\mathrm{Last} = 2943\,\mathrm{Nm}
\end{equation}
\noindent
Dabei gilt für $\Delta \omega$ als Differenz der Winkelgeschwindigkeiten vor und nach dem Bremsvorgang:
\begin{equation} \label{dw}
    |\Delta \omega| = |\omega_\mathrm{nach} - \omega_\mathrm{vor}| = |0 - \omega_\mathrm{ZW2}| = 95,771\,\mathrm{\frac{1}{s}} 
\end{equation}
Die Ergebnisse von (Gl.\ref{Jges}), (Gl. \ref{ML}) und (Gl. \ref{dw}) in (Gl. \ref{MK}) eingesetzt ergibt:
\begin{equation}
    M_\mathrm{K} = 40,372\,\mathrm{Nm}
\end{equation}





\subsection{Festlegung einiger Parameter}
Da es sich um eine nasslaufende Lamellenbremse handelt, muss zuerst eine passende Reibflächenpaarung ausgewählt werden. Aufgrund der höheren Gleitreibungszahl $\mu= 0.07$ gegenüber einer Stahl/Stahl-Paarung ($\mu = 0.05$) \cite[Tab. 1]{VDI2241-1}, bietet sich bei Nasslauf eine Stahl/Sinterbronze-Paarung an \cite[Tab. 1]{VDI2241-2}. \\
Es werden Innen- sowie Außenlamellen der Firma Ortlinghaus verwendet. Wegen des geringen verfügbaren Bauraums wird auf die kleinsten verfügbaren Lamellen zurückgegriffen; die wirksamen Innen- und Außendurchmesser ergeben sich demnach zu $D_\mathrm{a} = 54\,\mathrm{mm}$ und $d_\mathrm{i} = 34\,\mathrm{mm}$. \\
Für die weiteren Rechnungen wird zudem von einer Reibflächenzahl von $z = 12$ ausgegangen. 

\subsection{Normalkraft $F_\mathrm{N}$}
Durch die an den Lamellen wirkende Normalkraft $F_\mathrm{N}$ wird ein Reibmoment $M_\mathrm{R}$ hervorgerufen. Da dieses Reibmoment mindestens gleich dem Kennmoment $M_\mathrm{K}$ sein muss um eine Kupplung/Bremsung zu ermöglichen, lasst sich $F_\mathrm{N}$ aus dem Kennmoment herleiten \cite[S. 325, Gl. 5.42]{Sauer}:
\begin{equation}
    M_\mathrm{K} \stackrel{!}{=} M_\mathrm{R} = \mu \cdot r_\mathrm{m} \cdot z \cdot F_\mathrm{N}
\end{equation}
\begin{equation} \label{FN}
    F_\mathrm{N} = \frac{M_\mathrm{K}}{\mu \cdot r_\mathrm{m} \cdot z}
\end{equation}
Für den mittleren flächenbezogenen Reibflächendurchmesser $r_\mathrm{m}$ gilt vereinfacht \cite[S. 324]{Sauer}:
\begin{equation}
    r_\mathrm{m} = \frac{1}{2} \left( \frac{D_\mathrm{a}}{2} + \frac{d_\mathrm{i}}{2} \right) = 22\,\mathrm{mm}
\end{equation}
Somit folgt aus (Gl. \ref{FN}): 
\begin{equation}
    F_\mathrm{N} = 2184,648\,\mathrm{N}
\end{equation}


\subsection{Überprüfung der Einhaltung von Richtwerten}
In VDI 2241 \cite[Tab. 2]{VDI2241-2} sind einige Richtwerte für die verschiedenen Reibpaarung gegeben. Diese sollten nicht überschritten werden, um weitere Auslegungsrechnungen zu vermeiden. 
Folgende Richtwerte sind für eine Stahl/Sinterbronze-Paarung einzuhalten, d.h. nicht zu überschreiten \cite[S. 5, Tab. 2]{VDI2241-2}:
\begin{table}[H]
    \centering
    \caption{Richtwerte der Bremsauslegung}
    \begin{tabular}{p{8cm}ccc}
        \toprule
         Bezeichnung&  Formelzeichen& Größe& Einheit\\
         \midrule
         zulässige flächenbezogene Schaltarbeit (einmalige Schaltung)& $q_\mathrm{AE,zul}$& 1 bis 2& $[\mathrm{\frac{J}{mm^2}}]$\\
         max. Reibflächenpressung& $p_\mathrm{R,max}$& $4$& $[\mathrm{\frac{N}{mm^2}}]$\\
         max. Gleitgechwindigkeit& $v_\mathrm{R,max}$& $40$& $[\mathrm{\frac{m}{s}}]$ \\
         zulässige flächenbezogene Reibleistung& $\dot{q}_\mathrm{A,zul}$& 1,5 bis 2,5& $[\mathrm{\frac{W}{mm^2}}]$ \\
         \bottomrule
    \end{tabular}
    \label{Brumse_Richt}
\end{table}\\


\subsubsection{Flächenbezogene Schaltarbeit bei einmaliger Schaltung $q_\mathrm{AE}$} \label{Sec_qAE}
Aus den Gleichungen (4) und (8) aus VDI 2241 Blatt 1 \cite[S. 15]{VDI2241-1} ergibt sich nach trivialem Erweitern und Einsetzen eine vorhandene Schaltarbeit $Q_\mathrm{vorh}$: 
\begin{equation}
    Q_\mathrm{vorh} = \frac{1}{2} M_\mathrm{K} \cdot |\Delta \omega| \cdot t_3 =  1,927\,\mathrm{kJ}
\end{equation}
Die gesamte Reibfläche der Bremse ist nach ...:
\begin{equation}
    A_\mathrm{R} = \pi \cdot z \cdot \left( \left(\frac{D_\mathrm{a}}{2}\right)^2 - \left(\frac{d_\mathrm{i}}{2}\right)^2 \right) = 1,659\times10^{-2}\,\mathrm{m^2}
\end{equation} 
Die vorhandene flächenbezogene Schaltarbeit bei einmaliger Schaltung $q_\mathrm{AE}$:
\begin{equation}
    q_\mathrm{AE,vorh} = \frac{Q_\mathrm{vorh}}{A_\mathrm{R}} = 1,161\times 10^{-1}\,\mathrm{\frac{J}{mm^2}}
\end{equation}
Ein Vergleich mit Tab. \ref{Brumse_Richt} zeigt die Einhaltung des Richtwerts $q_\mathrm{AE,zul}$:
\begin{equation}
    q_\mathrm{AE,vorh} = 1,161\times10^{-1}\,\mathrm{\frac{J}{mm^2}} <  1\mathrm{\frac{J}{mm^2}} = q_\mathrm{AE,zul}
\end{equation}

\subsubsection{Reibflächenpressung $p_\mathrm{R}$}
Anders als in Kap. \ref{Sec_qAE} ist für die Reibflächenpressung lediglich die Fläche einer Lamellenpaarung zu berücksichtigen (vgl. \cite[S. 46]{RM}). Die Reibfläche $A$ einer einzelnen Lamellenpaarung ist:
\begin{equation}
    A = \pi \left( \left(\frac{D_\mathrm{a}}{2}\right)^2 - \left(\frac{d_\mathrm{i}}{2}\right)^2 \right) = 1,382\times10^{-3}\,\mathrm{m^2}
\end{equation}
Nach \cite[S. 46, Gl. 3.10]{RM} ergibt sich die vorhandene Reibflächenpressung zu:
\begin{equation}
    p_\mathrm{R,vorh} = \frac{F_\mathrm{N,vorh}}{A} = 1,580\,\mathrm{\frac{N}{mm^2}}
\end{equation}
Der Vergleich mit Tab. \ref{Brumse_Richt} zeigt:
\begin{equation}
    p_\mathrm{R,vorh} = 1,580\,\mathrm{\frac{N}{mm^2}} < 4\,\mathrm{\frac{N}{mm^2}} = p_\mathrm{R,max}
\end{equation}


\subsubsection{Gleitgeschwindigkeit $v_\mathrm{R}$}
Die Gleitgeschindwigkeit $v_\mathrm{R}$ berechnet sich als Produkt von mittlerem Reibradius $r_\mathrm{m}$ und der Winkelgeschwindigkeit $\omega_\mathrm{ZW2}$, mit der sich die Lamellen drehen:
\begin{equation}
    v_\mathrm{R,vorh} = r_\mathrm{m} \cdot \omega_\mathrm{ZW2} = 2,100\,\mathrm{\frac{m}{s}}
\end{equation}
Die maximale Gleitgeschwindigkeit wird also deutlich unterschritten:
\begin{equation}
    v_\mathrm{R,vorh} = 2,100\,\mathrm{\frac{m}{s}} < 40\,\mathrm{\frac{m}{s}} = v_\mathrm{R,max}
\end{equation}

\subsubsection{Flächenbezogene Reibleistung $\dot{q}_\mathrm{A}$}
Für die vorhandene flächenbezogene Reibleistung gilt nach \cite[S. 15, Gl. 12]{VDI2241-1}:
\begin{equation}
    \dot{q}_\mathrm{A,vorh} = v_\mathrm{R,vorh} \cdot p_\mathrm{R,vorh} \cdot \mu = 0,232\,\mathrm{\frac{W}{mm^2}}
\end{equation}
Auch dieser Richtwert wird eingehalten (vgl. Tab. \ref{Brumse_Richt}): 
\begin{equation}
    \dot{q}_\mathrm{A,vorh} =  0,232\,\mathrm{\frac{W}{mm^2}} < 2,5\,\mathrm{\frac{W}{mm^2}} = \dot{q}_\mathrm{A,zul}
\end{equation}\\
\newline
\noindent
Somit sind alle Richtwerte aus Tab. \ref{Brumse_Richt} (vgl. \cite[S. 5, Tab. 2]{VDI2241-2})eingehalten, und die Bremse entspricht in dieser Form der VDI 2241. 


\subsection{Auslegung der Druckfedern}
Die Auslegung der Druckfedern erfolgt über die Software KISSsoft; alle nachfolgenden Werte werden in diesem Programm nach DIN EN 16984:2017 berechnet. 
\newline
\subsubsection{Vorauslegung}
Vor der eigentlichen Auslegung der Federn, sind einige .... zu berechnen. 
Zuerst ist der notwendige Federweg $\Delta s$ zu bestimmen. Angenommen, der Anker der Bremse grenzt im gelüfteten Zustand unmittelbar an die äußerste Außenlamelle und drückt das Lamellenpaket im angezogenen Zustand maximal zusammen. Dann ist der notwendige $\Delta s$ gleich der Differenz der Breiten $\Delta B$ des Lamellenpakets in den beiden beschriebenen Zuständen. Folgende ... der Lamellen sind bekannt:
\begin{table}[H]
    \centering
    \caption{Lamellen}
    \begin{tabular}{p{8cm}ccc}
        \toprule
         Bezeichnung&  Formelzeichen& Größe& Einheit\\
         \midrule
         Anzahl wirksamer Reibflächen& $n$& $12$& \\
         Breite Außenlamellen& $b_\mathrm{aus}$& $1,5$& \\
         Breite Innenlamellen& $b_\mathrm{in}$& $1$& \\
         Sinushöhe Innenlamellen& $h_\mathrm{sin,in}$& $0,25$& \\
         Lüftspiel& $l_\mathrm{Luft}$& $0.2$& \\
         \bottomrule
    \end{tabular}
    \label{Lamellen}
\end{table}\\
\noindent
Die Gesamtbreite des Pakets im angezogenen Zustand berechnet sich demnach als:
\begin{equation}
   B_\mathrm{zu} = \frac{n}{2} \cdot b_\mathrm{in} + \left(\frac{n}{2} + 1\right) \cdot b_\mathrm{aus} = 16,5\,\mathrm{mm}
\end{equation}
Für die Lamellen im gelüfteten Zustand ergibt sich:
\begin{equation}
    B_\mathrm{auf} = \frac{n}{2} \cdot (b_\mathrm{in} + h_\mathrm{sin}) + \left(\frac{n}{2} + 1\right) \cdot b_\mathrm{aus} + l_\mathrm{Luft} = 18,5\,\mathrm{mm}
\end{equation}
Somit folgt für die Differenz der Breiten:
\begin{equation}
    \Delta B = \Delta s = B_\mathrm{auf} - B_\mathrm{zu} = 2\,\mathrm{mm}
\end{equation}


\subsubsection{Federwege und -kräfte}
Da sich die notwendige Normalkraft $F_\mathrm{N} = \,\mathrm{N}$ nur schwer genau erzielen lässt, erfolgt die Auslegung der Federn über die Federwege. Durch den kleinen Federweg $s_1$ muss dabei mindestens die Normalkraft aufgebracht werden. Für den großen Federweg $s_2$ folgt also:
\begin{equation}
    s_2 = s_1 + \Delta s
\end{equation}

\noindent
Folgende Federkennwerte sind Eingabewerte der Berechnung in KISSsoft:  
\begin{table}[H]
    \centering
    \caption{festgelegte Federkennwerte}
    \begin{tabular}{p{8cm}ccc}
        \toprule
        Bezeichnung& Formelzeichen& Größe& Bemerkung \\
        \midrule
        Außendurchmesser& $D_\mathrm{e}$& $71$& \\
        Innendurchmesser& $D_\mathrm{i}$& $36$& \\
        Vergrößerter Innendurchmesser& $D_\mathrm{i'}$&  $40$& \\
        Dicke Einzelteller& $t$& $2,5$& \\
        Höher Einzelteller& $l_0$& $4,5$& \\
        Federn pro Paket& $n$& $1$& \\
        Federpakete pro Säule& $i$& $3$& \\
        \midrule
        kleiner Federweg& $s_1$& $2,2$& \\
        großer Federweg& $s_2$& $4,2$& \\
        \bottomrule
    \end{tabular}
    \label{KISSsoft Federn Eingabe}
\end{table}\\


\noindent
Die Ergebnisse dieser Berechnungen sind:
\begin{table}[H]
    \centering
    \caption{Errechnete Federkennwerte}
    \begin{tabular}{p{8cm}ccc}
        \toprule
        Bezeichnung&  Formelzeichen& Größe& Bemerkung\\
        \midrule
         kleine Federkraft& $F_1$& $2581,856\,\mathrm{N}$& \\
         großer Federweg& $F_2$& $4295,974\,\mathrm{N}$& \\
        \midrule
         Gesamtlänge unbelastet& $L_0$& $13,5\,\mathrm{mm}$& \\
         Gesamtlänge nach $F_1$& $L_1$& $11,3\,\mathrm{mm}$& \\
         Gesamtlänge nach $F_2$& $L_2$& $9,3\,\mathrm{mm}$& \\
        \bottomrule
    \end{tabular}
    \label{KISSsoft Federn Ausgabe}
\end{table}\\

\noindent
Aus Tab. (\ref{KISSsoft Federn Ausgabe}) geht hervor, dass 
\begin{equation*}
    F_1 > F_\mathrm{N}
\end{equation*}
und somit genug Anpresskraft für eine Bremsung vorhanden ist.

\subsection{Hydraulik}
Um die Bremse vollständig lüften zu können, muss die durch den hydraulischen Druck hervorgerufene Kraft $F_\mathrm{hyd}$ größer sein als die große Federkraft $F_2$. 
Für die Druckkraft gilt nach \cite[S. 46, Gl. 3.10]{RM}:
\begin{equation} \label{Fhyd}
    F_\mathrm{hyd} = p_\mathrm{hyd} \cdot A_\mathrm{hyd}
\end{equation}
Für den minimal erforderlichen Druck folgt der Zusammenhang: 
\begin{equation} \label{phyd}
    p_\mathrm{hyd,min} \geq \frac{F_2}{A_\mathrm{hyd}}
\end{equation}
Die wirksame Druckfläche $A_\mathrm{hyd}$ der Hydraulikkammer ergibt sich aus der Geometrie des Ankers:
\begin{equation}
    A_\mathrm{hyd} = \pi \left( \left(\frac{D_\mathrm{Anker}}{2} \right)^2 - \left(\frac{d_\mathrm{Anker}}{2}\right)^2 \right)
\end{equation}
Die Geometrie des Ankers folgt aus der Modellierung in Solid Edge:
\begin{table}[H]
    \centering
    \caption{Ankergeometrie}
    \begin{tabular}{p{8cm}ccc}
        \toprule
        Bezeichnung&  Formelzeichen& Größe& Bemerkung\\
        \midrule
         großer Außendruchmesser& $D_\mathrm{Anker}$& $80\,\mathrm{mm}$& \\
         kleiner Außendurchmesser& $d_\mathrm{Anker}$& $54\,\mathrm{mm}$& \\
        \bottomrule
    \end{tabular}
    \label{Ankergeometrie}
\end{table}\\

\noindent
Für die Druckfläche folgt demnach: 
\begin{equation}
    A_\mathrm{hyd} = 2736\,\mathrm{mm^2}
\end{equation}

\noindent
Der minimale hydraulische Druck ist nach Gl. (\ref{phyd}):
\begin{equation}
    p_\mathrm{hyd,min} \geq 1,57\,\mathrm{\frac{N}{mm^2}} = 15,7\,\mathrm{bar}
\end{equation}

\noindent
Es liegt also nah, einen tatsächlichen hydraulischen Druck von $p_\mathrm{hyd} = 16\,\mathrm{bar}$ zu wählen.
Durch Einsetzen in Gl. (\ref{Fhyd}) ergibt sich die vorhandene Druckkraft:
\begin{equation}
    F_\mathrm{hyd,vorh} = 4378\,\mathrm{N}
\end{equation}
Somit ist also auch die Bedingung 
\begin{equation*}
    F_\mathrm{hyd,vorh} \geq F_2
\end{equation*}
erfüllt (vgl. Tab. (\ref{KISSsoft Federn Ausgabe})), und die Federn können vollständig vorgespannt werden. 